\textbf{Phần 1. Trắc nghiệm lựa chọn}
\setcounterpageref{ex}{0}

\begin{ex} %[3H3V43]
Một đoạn thẳng $MN$ dài $64\text{ cm}$. Điểm $I$ là trung điểm của $MN$, điểm $K$ là trung điểm của $MI$. Độ dài đoạn thẳng $KN$ là
\choice
{$32\text{ cm}$}
{$40\text{ cm}$}
{\True $48\text{ cm}$}
{$56\text{ cm}$}
\loigiai{
$MI = 64 : 2 = 32\text{ cm}$; $MK = 32 : 2 = 16\text{ cm}$; $KN = MN - MK = 64 - 16 = 48\text{ (cm)}$.
\textbf{Đáp án của câu hỏi là C}
}
\end{ex}

\begin{ex} %[3H3V43]
Một hình tròn có đường kính $18\text{ cm}$. Bán kính của hình tròn đó là
\choice
{\True $9\text{ cm}$}
{$36\text{ cm}$}
{$6\text{ cm}$}
{$12\text{ cm}$}
\loigiai{
$18 : 2 = 9\text{ (cm)}$.
\textbf{Đáp án của câu hỏi là A}
}
\end{ex}

\begin{ex} %[3H3V43]
Một hình khối lập phương có bao nhiêu đỉnh và bao nhiêu cạnh?
\choice
{$6$ đỉnh và $12$ cạnh}
{$8$ đỉnh và $6$ cạnh}
{$12$ đỉnh và $8$ cạnh}
{\True $8$ đỉnh và $12$ cạnh}
\loigiai{
Khối lập phương có 8 đỉnh và 12 cạnh.
\textbf{Đáp án của câu hỏi là D}
}
\end{ex}

\begin{ex} %[3D3V43]
Đổi $2\text{ m } 4\text{ cm}$ ra đơn vị mi-li-mét:
\choice
{$240\text{ mm}$}
{\True $2040\text{ mm}$}
{$20400\text{ mm}$}
{$2400\text{ mm}$}
\loigiai{
$2\text{ m} = 2000\text{ mm}$; $4\text{ cm} = 40\text{ mm}$. Tổng: $2040\text{ mm}$.
\textbf{Đáp án của câu hỏi là B}
}
\end{ex}

\begin{ex} %[3D3V43]
Một gói đường nặng $1\text{ kg}$. Người ta lấy ra $350\text{ g}$ nấu chè và $250\text{ g}$ làm bánh. Gói đường còn lại bao nhiêu gam?
\choice
{$500\text{ g}$}
{$300\text{ g}$}
{\True $400\text{ g}$}
{$600\text{ g}$}
\loigiai{
Đổi $1\text{ kg} = 1000\text{ g}$. Đã dùng: $350 + 250 = 600\text{ g}$. Còn lại: $1000 - 600 = 400\text{ (g)}$.
\textbf{Đáp án của câu hỏi là C}
}
\end{ex}

\begin{ex} %[3D3V43]
Một bình nước có dung tích $1\ell$. Rót đầy 4 cốc nước loại $200\text{ ml}$. Trong bình còn lại bao nhiêu mi-li-lít nước?
\choice
{\True $200\text{ ml}$}
{$100\text{ ml}$}
{$300\text{ ml}$}
{$400\text{ ml}$}
\loigiai{
Đổi $1\ell = 1000\text{ ml}$. Bốn cốc chứa: $200 \times 4 = 800\text{ ml}$. Còn lại: $1000 - 800 = 200\text{ (ml)}$.
\textbf{Đáp án của câu hỏi là A}
}
\end{ex}

\begin{ex} %[3D3V43]
Nhiệt độ nước đá đang tan là $0^\circ\text{C}$, nhiệt độ nước sôi là $100^\circ\text{C}$. Mức chênh lệch nhiệt độ giữa nước sôi và nước đá đang tan là
\choice
{$50^\circ\text{C}$}
{$0^\circ\text{C}$}
{$90^\circ\text{C}$}
{\True $100^\circ\text{C}$}
\loigiai{
$100 - 0 = 100^\circ\text{C}$.
\textbf{Đáp án của câu hỏi là D}
}
\end{ex}

\begin{ex} %[3H3V43]
Một hình chữ nhật có chu vi $50\text{ cm}$, chiều dài là $15\text{ cm}$. Chiều rộng của hình chữ nhật là
\choice
{$20\text{ cm}$}
{\True $10\text{ cm}$}
{$5\text{ cm}$}
{$12\text{ cm}$}
\loigiai{
Nửa chu vi: $50 : 2 = 25\text{ cm}$. Chiều rộng: $25 - 15 = 10\text{ (cm)}$.
\textbf{Đáp án của câu hỏi là B}
}
\end{ex}

\begin{ex} %[3H3V43]
Góc tạo bởi kim giờ và kim phút lúc đúng 3 giờ là
\choice
{Góc nhọn}
{Góc tù}
{\True Góc vuông}
{Góc bẹt}
\loigiai{
Lúc 3 giờ đúng hai kim vuông góc với nhau.
\textbf{Đáp án của câu hỏi là C}
}
\end{ex}

\begin{ex} %[3H3V43]
Một khối gỗ hình lập phương có cạnh $6\text{ cm}$. Tổng độ dài tất cả các cạnh của khối gỗ là
\choice
{\True $72\text{ cm}$}
{$36\text{ cm}$}
{$48\text{ cm}$}
{$24\text{ cm}$}
\loigiai{
Khối lập phương có 12 cạnh: $12 \times 6 = 72\text{ (cm)}$.
\textbf{Đáp án của câu hỏi là A}
}
\end{ex}

\begin{ex} %[3D3V43]
Một con gà cân nặng $2\text{ kg } 300\text{ g}$, con ngỗng cân nặng $4\text{ kg } 100\text{ g}$. Cả gà và ngỗng cân nặng là
\choice
{$6\text{ kg } 200\text{ g}$}
{$6\text{ kg } 500\text{ g}$}
{$7\text{ kg}$}
{\True $6\text{ kg } 400\text{ g}$}
\loigiai{
$2300 + 4100 = 6400\text{ g} = 6\text{ kg } 400\text{ g}$.
\textbf{Đáp án của câu hỏi là D}
}
\end{ex}

\begin{ex} %[3H3V43]
Người ta uốn một đoạn dây thép dài $60\text{ cm}$ thành một hình tam giác có 3 cạnh bằng nhau. Độ dài mỗi cạnh tam giác là
\choice
{$15\text{ cm}$}
{\True $20\text{ cm}$}
{$30\text{ cm}$}
{$25\text{ cm}$}
\loigiai{
$60 : 3 = 20\text{ (cm)}$.
\textbf{Đáp án của câu hỏi là B}
}
\end{ex}

\textbf{Phần 2. Trắc nghiệm đúng sai}
\setcounterpageref{ex}{0}

\begin{ex} %[3H3V43]
Xét tính đúng sai của các kiến thức hình học và đo lường:
\choiceTF
{\True Hình vuông có 4 cạnh bằng nhau và 4 góc vuông.}
{Đoạn thẳng nối hai điểm bất kì trên đường tròn là đường kính.}
{\True $1\text{ kg} = 1000\text{ g}$ và $1\ell = 1000\text{ ml}$.}
{$1\text{ m} = 100\text{ mm}$.}
\loigiai{
\begin{itemize}
    \item Ý a) Đúng: Định nghĩa hình vuông.
    \item Ý b) Sai: Sai vì chỉ đoạn qua tâm mới là đường kính.
    \item Ý c) Đúng: Quan hệ chuẩn.
    \item Ý d) Sai: Sai vì $1\text{ m} = 1000\text{ mm}$.
\end{itemize}
\begin{itemize}
    \item a) \textbf{Đúng}.
    \item b) \textbf{Sai}.
    \item c) \textbf{Đúng}.
    \item d) \textbf{Sai}.
\end{itemize}
}
\end{ex}

\begin{ex} %[3H3V43]
Một hình chữ nhật có chiều dài $18\text{ cm}$, chiều rộng bằng $\dfrac{1}{3}$ chiều dài:
\choiceTF
{\True Chiều rộng của hình chữ nhật bằng $6\text{ cm}$.}
{\True Chu vi của hình chữ nhật bằng $48\text{ cm}$.}
{Có thể chia hình chữ nhật đó thành 2 hình vuông cạnh $6\text{ cm}$.}
{\True Nửa chu vi của hình chữ nhật bằng $24\text{ cm}$.}
\loigiai{
\begin{itemize}
    \item Ý a) Đúng: $18 : 3 = 6\text{ cm}$.
    \item Ý b) Đúng: $(18 + 6) \times 2 = 48\text{ cm}$.
    \item Ý c) Sai: Sai vì $18 : 6 = 3$ hình vuông.
    \item Ý d) Đúng: $48 : 2 = 24\text{ cm}$.
\end{itemize}
\begin{itemize}
    \item a) \textbf{Đúng}.
    \item b) \textbf{Đúng}.
    \item c) \textbf{Sai}.
    \item d) \textbf{Đúng}.
\end{itemize}
}
\end{ex}

\textbf{Phần 3. Trả lời ngắn}
\setcounterpageref{ex}{0}

\begin{ex} %[3H3V43]
Một hình vuông có chu vi $56\text{ cm}$. Tính độ dài cạnh hình vuông theo xăng-ti-mét.
\shortans{14}
\loigiai{
\textbf{Đáp số:} $14$\\
$56 : 4 = 14\text{ (cm)}$.
}
\end{ex}

\begin{ex} %[3H3V43]
Hình tròn có bán kính $15\text{ cm}$. Hỏi đường kính dài bao nhiêu xăng-ti-mét?
\shortans{30}
\loigiai{
\textbf{Đáp số:} $30$\\
$15 \times 2 = 30\text{ (cm)}$.
}
\end{ex}

\begin{ex} %[3D3V43]
Đổi $1\text{ m } 8\text{ cm}$ ra đơn vị mi-li-mét.
\shortans{1080}
\loigiai{
\textbf{Đáp số:} $1080$\\
$1000 + 80 = 1080\text{ mm}$.
}
\end{ex}

\begin{ex} %[3D3V43]
Một túi gạo nặng $2\text{ kg}$. Người ta đã dùng $750\text{ g}$. Hỏi túi gạo còn lại bao nhiêu gam?
\shortans{1250}
\loigiai{
\textbf{Đáp số:} $1250$\\
$2000 - 750 = 1250\text{ (g)}$.
}
\end{ex}

\begin{ex} %[3D3V43]
Một bình chứa $1\ell\text{ }200\text{ ml}$ nước. Rót đều vào 4 cốc. Mỗi cốc chứa bao nhiêu mi-li-lít nước?
\shortans{300}
\loigiai{
\textbf{Đáp số:} $300$\\
$1200 : 4 = 300\text{ (ml)}$.
}
\end{ex}

\begin{ex} %[3H3V43]
Một khối lập phương có chu vi một mặt là $24\text{ cm}$. Tính tổng độ dài 12 cạnh của khối lập phương đó theo xăng-ti-mét.
\shortans{72}
\loigiai{
\textbf{Đáp số:} $72$\\
Cạnh: $24 : 4 = 6\text{ cm}$. Tổng 12 cạnh: $12 \times 6 = 72\text{ (cm)}$.
}
\end{ex}
